\chapter{Matrices as Machines: Linear Transformations}

So far our vectors have been static arrows resting in space. You have added
them, scaled them, combined them into spans, and used the dot product to measure
their lengths, angles, and projections. But throughout all of that, the space
itself sat still. Now we ask the question that turns linear algebra from static
geometry into a living, moving world: \emph{what if space itself can move? What
if a matrix is not a static grid of numbers, but a \textbf{verb}?}

A matrix represents a transformation. You supply an input vector, and it produces
an output vector. Apply it to $\colvec{1\\0}$ and see where it lands. Apply it to
$\colvec{0\\1}$ and observe the result. Extending this across the entire plane,
the transformation moves all of space at once: rotating, stretching, reflecting, or
collapsing it. Everything else in this book (elimination, inverses, eigenvectors, the
SVD) builds on understanding what these linear transformations do and how to decompose them.

This chapter builds that picture from the ground up. We start with the geometry
of a single transformation, see why linearity preserves grids, and examine
the basic motions of the plane: rotations, reflections, shears, and scalings.

\section{A Matrix Is a Function on Vectors}
\label{ch05:matrix-is-function}

Here is the whole idea in one example. Take the matrix
\[
  \mat{A} = \begin{bmatrix} 2 & 0 \\ 0 & 3 \end{bmatrix},
\]
and feed it the vector $\vec{x} = \colvec{1\\1}$. How do we feed a vector into
a matrix? Each row of $\mat{A}$ takes a dot product with $\vec{x}$ -- the very
operation we mastered in the previous chapter:
\[
  \mat{A}\vec{x}
  = \begin{bmatrix} 2 & 0 \\ 0 & 3 \end{bmatrix}\colvec{1\\1}
  = \colvec{\begin{bmatrix} 2 & 0 \end{bmatrix} \cdot \colvec{1\\1} \\ \begin{bmatrix} 0 & 3 \end{bmatrix} \cdot \colvec{1\\1}}
  = \colvec{2\cdot 1 + 0 \cdot 1 \\ 0 \cdot 1 + 3 \cdot 1}
  = \colvec{2\\3}.
\]
Each row computes a dot product with $\vec{x}$ to produce one coordinate of the
output. The matrix took the arrow pointing to $(1,1)$ and returned the arrow
pointing to $(2,3)$: it doubled the horizontal coordinate and tripled the vertical
coordinate. Applied to any other vector, it scales the coordinates in the same way.
The matrix is a rule that assigns to \emph{every} input vector $\vec{x}$ an output
vector $\mat{A}\vec{x}$. That is precisely what a function is. In geometry, a function
whose inputs and outputs are vectors is called a \textbf{transformation} -- it transforms
space by mapping every vector to an output vector.

\begin{intuitionbox}
Think of $\mat{A}\vec{x}$ geometrically as \textbf{$\mat{A}$ acting on $\vec{x}$}
-- the transformation $\mat{A}$ mapping the point $\vec{x}$ to its destination
$\mat{A}\vec{x}$. The matrix describes the overall transformation of space; the
vector $\vec{x}$ is one point being mapped.
\end{intuitionbox}

We will write this as a function and give it a name. An $m \times n$ matrix
$\mat{A}$ defines a function
\[
  T : \R^n \to \R^m, \qquad T(\vec{x}) = \mat{A}\vec{x}.
\]
It takes in vectors with $n$ entries and produces vectors with $m$ entries -- the
column count $n$ is the dimension of the input space, while the row count $m$ is
the dimension of the output space. A $2\times 2$ matrix maps the plane to the plane.
A $3\times 2$ matrix maps the plane into $3$-dimensional space. A $1\times 3$
matrix maps $3$-space to a single real number (a line). The \emph{dimensions} of
the matrix tell you which spaces it connects before you compute a single entry.

\begin{focusbox}
$\mat{A}$ is $m \times n$ $\;\Longrightarrow\;$ input lives in $\R^n$ (the
columns), output lives in $\R^m$ (the rows). Read the shape first: it tells you
where vectors come from and where they go.
\end{focusbox}

\begin{example}[Three machines, one vector]
\label{ch05:three-machines}
Let $\vec{x} = \colvec{2\\1}$ and watch three different matrices act on it.

\emph{A flip across the vertical axis,} $\mat{F} = \begin{bmatrix} -1 & 0 \\ 0 & 1\end{bmatrix}$:
\[
  \mat{F}\vec{x} = \colvec{-1\cdot 2 + 0 \\ 0 + 1\cdot 1} = \colvec{-2\\1}.
\]
The $x$-coordinate changed sign; the point hopped to the other side of the
$y$-axis.

\emph{A horizontal shear,} $\mat{S} = \begin{bmatrix} 1 & 1 \\ 0 & 1\end{bmatrix}$:
\[
  \mat{S}\vec{x} = \colvec{1\cdot 2 + 1 \cdot 1 \\ 0 + 1 \cdot 1} = \colvec{3\\1}.
\]
The height $1$ slid the point to the right by $1$.

\emph{A collapse onto the $x$-axis,} $\mat{P} = \begin{bmatrix} 1 & 0 \\ 0 & 0\end{bmatrix}$:
\[
  \mat{P}\vec{x} = \colvec{2\\0}.
\]
The vertical part was erased; every point now lives on the horizontal line.
The same input vector lands in three different places depending on which matrix
acts on it.
\end{example}

\subsection{The picture: dots moving to dots}

The cleanest way to \emph{see} a transformation is to draw a few input points and
draw arrows to where they land. Figure~\ref{ch05:fig-action} shows the doubling
machine $\mat{A} = \diag(2,3)$ acting on a small cloud of points.

\begin{figure}[H]
\centering
\begin{tikzpicture}[scale=0.9,>=Stealth]
  \begin{scope}
    \draw[->,gray!60] (-0.4,0) -- (5.2,0) node[right]{$x$};
    \draw[->,gray!60] (0,-0.4) -- (0,6.6) node[above]{$y$};
    % input points (blue) and output points (green) with arrows
    \foreach \px/\py in {1/0, 0/1, 1/1, 2/1, 1/2} {
      \pgfmathsetmacro{\qx}{2*\px}
      \pgfmathsetmacro{\qy}{3*\py}
      \fill[vecblue] (\px,\py) circle (2pt);
      \fill[targetgreen] (\qx,\qy) circle (2pt);
      \draw[->,gray!50,dashed] (\px,\py) -- (\qx,\qy);
    }
    \node[vecblue,below left] at (1,1) {$(1,1)$};
    \node[targetgreen,above right] at (2,3) {$(2,3)$};
  \end{scope}
\end{tikzpicture}
\caption{The machine $\mat{A}=\protect\begin{bmatrix}2&0\\0&3\protect\end{bmatrix}$
acting on a handful of points. Blue dots are inputs, green dots are outputs, and
each dashed arrow shows where a point is dragged. Horizontal distances double,
vertical distances triple.}
\label{ch05:fig-action}
\end{figure}

A cloud of dots tells the truth but it gets cluttered fast. The trick that makes
transformations \emph{readable} is to draw a whole \textbf{grid} and watch what
the machine does to it. Because (as we will prove in a moment) the machine moves straight lines to
straight lines and evenly spaced lines to evenly spaced lines, the image of a
square grid is again a grid -- but possibly tilted, stretched, or sheared. We
will lean on this picture constantly.

\section{Linearity: The One Property That Makes a Matrix a Matrix}
\label{ch05:linearity}

Not every function on vectors comes from a matrix. The function that adds
$\colvec{1\\1}$ to everything, or the function that maps $\vec{x}$ to the vector
of squared entries, are perfectly good functions -- but they are not matrices.
The functions that \emph{are} matrices are exactly the ones with a special,
beautiful property called \textbf{linearity}.

\begin{definition}[Linear transformation]
\label{ch05:def-linear}
A function $T : \R^n \to \R^m$ is \textbf{linear} if for all vectors
$\vec{u}, \vec{v} \in \R^n$ and all scalars $c \in \R$:
\begin{align*}
  \textbf{(additivity)} \qquad & T(\vec{u} + \vec{v}) = T(\vec{u}) + T(\vec{v}), \\
  \textbf{(homogeneity)} \qquad & T(c\,\vec{u}) = c\,T(\vec{u}).
\end{align*}
\end{definition}

Read these geometrically. \emph{Additivity} says: if you build $\vec{u}+\vec{v}$
by the parallelogram rule and \emph{then} transform, you get the same arrow as
transforming $\vec{u}$ and $\vec{v}$ separately and \emph{then} adding. The
machine doesn't care whether you combine before or after (Figure~\ref{ch05:fig-additivity}).
\emph{Homogeneity} says: stretching the input by $c$ stretches the output by the
same $c$. Double the input arrow, double the output arrow. The machine respects
scaling.

\begin{figure}[H]
\centering
\begin{tikzpicture}[scale=0.9,>=Stealth]
  % Input plane
  \begin{scope}[shift={(0,0)}]
    \draw[->,gray!50] (-0.3,0) -- (3.2,0) node[right]{$x$};
    \draw[->,gray!50] (0,-0.3) -- (0,2.6) node[above]{$y$};
    \draw[gray!40,dashed] (1.6,0.4) -- (2.2,1.8) -- (0.6,1.4);
    \draw[->,vecblue,very thick] (0,0) -- (1.6,0.4) node[below right]{$\vec{u}$};
    \draw[->,vecred,very thick] (0,0) -- (0.6,1.4) node[above left]{$\vec{v}$};
    \draw[->,quotientorange,very thick] (0,0) -- (2.2,1.8) node[above right]{$\vec{u}+\vec{v}$};
    \node[font=\small] at (1.5,-0.8) {add, then transform};
  \end{scope}

  % Arrow T
  \draw[->,thick,gray] (3.8,1.1) -- (5.2,1.1) node[midway,above]{$T$};

  % Output plane
  \begin{scope}[shift={(6.2,0)}]
    \draw[->,gray!50] (-0.3,0) -- (3.4,0) node[right]{$x$};
    \draw[->,gray!50] (0,-0.3) -- (0,2.6) node[above]{$y$};
    \draw[gray!40,dashed] (1.8,0.2) -- (2.6,2.0) -- (0.8,1.8);
    \draw[->,vecblue,very thick] (0,0) -- (1.8,0.2) node[below right]{$T(\vec{u})$};
    \draw[->,vecred,very thick] (0,0) -- (0.8,1.8) node[above left]{$T(\vec{v})$};
    \draw[->,quotientorange,very thick] (0,0) -- (2.6,2.0) node[above right]{$T(\vec{u}+\vec{v})$};
    \node[font=\small] at (1.6,-0.8) {transform, then add};
  \end{scope}
\end{tikzpicture}
\caption{The geometric meaning of additivity: the parallelogram rule survives the transformation.
Adding $\vec{u}$ and $\vec{v}$ first and then applying $T$ reaches the exact same arrow as transforming
$\vec{u}$ and $\vec{v}$ individually and adding their images: $T(\vec{u}+\vec{v}) = T(\vec{u}) + T(\vec{v})$.}
\label{ch05:fig-additivity}
\end{figure}

\begin{intuitionbox}
Linearity is the statement that \textbf{the grid stays a grid}. Straight,
evenly spaced lines go to straight, evenly spaced lines, and the origin stays
put. No bending, no curving, no off-center shift. If you can picture a sheet of
graph paper being stretched and sheared but never crumpled, you are picturing a
linear transformation.
\end{intuitionbox}

The two rules combine into one clean statement worth memorizing: a linear map
sends linear combinations to linear combinations,
\[
  T(c\,\vec{u} + d\,\vec{v}) = c\,T(\vec{u}) + d\,T(\vec{v}).
\]
This is the heart of everything. It says $T$ \emph{commutes} with the basic
operations of linear algebra -- adding and scaling -- which are the only
operations we have. To say two operations \textbf{commute} means doing them
in either order yields the identical result: you can combine first and then
transform, or transform first and then combine. A linear map is a function
that preserves vector addition and scalar multiplication.

\begin{trapbox}
``Linear'' here does \emph{not} mean ``graph is a line.'' The high-school line
$y = mx + b$ is \textbf{not} a linear map when $b \neq 0$, because it fails
homogeneity: $T(0) = b \neq 0$, but every linear map must send $\vec{0}$ to
$\vec{0}$ (set $c = 0$ in homogeneity). An \textbf{affine map} is a linear
transformation followed by a rigid shift (a translation): it keeps straight lines
straight, but it drags the origin away from $(0,0)$. Linear maps always pin the
origin in place.
\end{trapbox}

\begin{example}[Checking linearity directly]
\label{ch05:check-linear}
Is rotation by $90^\circ$ counterclockwise a linear map? Call it $R$. Geometrically
$R$ takes $\colvec{a\\b}$ to $\colvec{-b\\a}$ (we will derive this carefully in
Section~\ref{ch05:rotations}). Check additivity:
\[
  R\!\left(\colvec{a_1\\b_1} + \colvec{a_2\\b_2}\right)
  = R\colvec{a_1+a_2\\b_1+b_2}
  = \colvec{-(b_1+b_2)\\a_1+a_2}
  = \colvec{-b_1\\a_1} + \colvec{-b_2\\a_2}
  = R\colvec{a_1\\b_1} + R\colvec{a_2\\b_2}.\ \checkmark
\]
Check homogeneity:
\[
  R\!\left(c\colvec{a\\b}\right) = R\colvec{ca\\cb} = \colvec{-cb\\ca}
  = c\colvec{-b\\a} = c\,R\colvec{a\\b}. \ \checkmark
\]
Both hold, so rotation is linear. Geometrically this is obvious: spinning the
whole plane rigidly about the origin keeps the grid a grid.
\end{example}

\begin{example}[A function that is \emph{not} linear]
\label{ch05:not-linear}
Let $T\colvec{a\\b} = \colvec{a^2\\b}$. Try homogeneity with $c = 2$ and
$\vec{u} = \colvec{1\\0}$:
\[
  T(2\vec{u}) = T\colvec{2\\0} = \colvec{4\\0},
  \qquad\text{but}\qquad
  2\,T(\vec{u}) = 2\colvec{1\\0} = \colvec{2\\0}.
\]
Since $\colvec{4\\0} \neq \colvec{2\\0}$, the map fails homogeneity. Squaring
bends the grid, so $T$ is not linear and cannot be written as a matrix.
\end{example}

\subsection{Why matrices are automatically linear}

Every matrix gives a linear map, and this property falls right out of
how matrix--vector products work. Two facts about the product $\mat{A}\vec{x}$ --
which we can take as the defining algebra of matrices -- are
\[
  \mat{A}(\vec{u} + \vec{v}) = \mat{A}\vec{u} + \mat{A}\vec{v}
  \qquad\text{and}\qquad
  \mat{A}(c\,\vec{u}) = c\,(\mat{A}\vec{u}).
\]
Compare with Definition~\ref{ch05:def-linear}: these are \emph{exactly}
additivity and homogeneity. So $T(\vec{x}) = \mat{A}\vec{x}$ is always linear.
Let us verify the first one by direct computation in $2\times 2$, so you trust the
machinery. Write $\mat{A} = \begin{bmatrix} a & b \\ c & d \end{bmatrix}$,
$\vec{u} = \colvec{u_1\\u_2}$, $\vec{v} = \colvec{v_1\\v_2}$:
\[
  \mat{A}(\vec{u}+\vec{v})
  = \begin{bmatrix} a & b \\ c & d \end{bmatrix}\colvec{u_1+v_1\\u_2+v_2}
  = \colvec{a(u_1+v_1) + b(u_2+v_2) \\ c(u_1+v_1) + d(u_2+v_2)}.
\]
Distribute and regroup by $\vec{u}$ and $\vec{v}$:
\[
  = \colvec{(au_1 + bu_2) + (av_1 + bv_2) \\ (cu_1+du_2) + (cv_1+dv_2)}
  = \colvec{au_1+bu_2\\cu_1+du_2} + \colvec{av_1+bv_2\\cv_1+dv_2}
  = \mat{A}\vec{u} + \mat{A}\vec{v}. \ \checkmark
\]
The only thing we used was the distributive law for ordinary numbers. Homogeneity
follows just as quickly by pulling the scalar out of each term:
\[
  \mat{A}(c\,\vec{u})
  = \begin{bmatrix} a & b \\ c & d \end{bmatrix}\colvec{c u_1 \\ c u_2}
  = \colvec{a(c u_1) + b(c u_2) \\ c(c u_1) + d(c u_2)}
  = c\colvec{au_1+bu_2\\cu_1+du_2}
  = c\,(\mat{A}\vec{u}). \ \checkmark
\]
Matrices are linear because arithmetic is distributive.

The centerpiece of the chapter is that the converse holds too: \emph{every}
linear map from $\R^n$ to $\R^m$ is given by some matrix. There are no other
linear maps hiding out there. To see why, and to learn how to \emph{build} the
matrix, we read the columns.

\section{What the Columns Tell You}
\label{ch05:columns}

Here is the single most useful sentence in this chapter:

\begin{focusbox}
\textbf{The columns of a matrix are the images of the standard basis vectors.}
Column $j$ of $\mat{A}$ is exactly $\mat{A}\vec{e}_j$ -- where the input arrow
$\vec{e}_j$ lands. To know a linear map, you only need to know where it sends the
basis vectors. Everything else follows by linearity.
\end{focusbox}

Recall the \textbf{standard basis vectors} of $\R^2$ -- the unit-length measuring
sticks pointing along each coordinate axis: $\vec{e}_1 = \colvec{1\\0}$ (one
step east) and $\vec{e}_2 = \colvec{0\\1}$ (one step north). Multiply a matrix by
$\vec{e}_1$ and watch what survives:
\[
  \begin{bmatrix} a & b \\ c & d \end{bmatrix}\colvec{1\\0}
  = \colvec{a\cdot 1 + b \cdot 0 \\ c \cdot 1 + d \cdot 0}
  = \colvec{a\\c} = \text{first column.}
\]
Likewise $\mat{A}\vec{e}_2 = \colvec{b\\d}$ is the second column. Picking out
$\vec{e}_j$ simply \emph{reads off} column $j$. This is why we say the columns
are the images of the basis vectors: feeding the machine the pure ``east'' arrow
returns its first column; feeding it the pure ``north'' arrow returns its second.

\begin{recoverybox}
\textbf{Mental check:} If a machine $\mat{A}$ sends $\vec{e}_1$ to $\colvec{5\\0}$
and $\vec{e}_2$ to $\colvec{0\\-2}$, what is the matrix $\mat{A}$?
(Pause before reading on: just drop the two destination arrows in as the two
columns: $\mat{A} = \begin{bmatrix} 5 & 0 \\ 0 & -2 \end{bmatrix}$.)
\end{recoverybox}

\subsection{From basis images to the whole map}

Now watch how a single matrix's behavior on the two basis vectors determines its
behavior on \emph{every} vector. Any $\vec{x} = \colvec{x_1\\x_2}$ is a linear
combination of the basis vectors:
\[
  \vec{x} = \colvec{x_1\\x_2} = x_1\colvec{1\\0} + x_2\colvec{0\\1}
          = x_1\vec{e}_1 + x_2\vec{e}_2.
\]
Apply $T$ and use linearity (Definition~\ref{ch05:def-linear}) to slide $T$
through the combination:
\[
  T(\vec{x}) = T(x_1\vec{e}_1 + x_2\vec{e}_2)
             = x_1\,T(\vec{e}_1) + x_2\,T(\vec{e}_2).
\]
Read that last expression aloud: \emph{the output is the same combination of the
basis images}. The numbers $x_1, x_2$ that built the input out of $\vec{e}_1,
\vec{e}_2$ now rebuild the output out of $T(\vec{e}_1), T(\vec{e}_2)$. So if you
stack the basis images as columns into a matrix $\mat{A} =
[\,T(\vec{e}_1)\ \ T(\vec{e}_2)\,]$, then
\[
  T(\vec{x}) = x_1\,(\text{col }1) + x_2\,(\text{col }2) = \mat{A}\vec{x}.
\]

\begin{formalbox}
\textbf{Every linear map is a matrix.} Let $T : \R^n \to \R^m$ be linear. Define
the $m \times n$ matrix $\mat{A}$ whose $j$-th column is $T(\vec{e}_j)$. Then
$T(\vec{x}) = \mat{A}\vec{x}$ for all $\vec{x} \in \R^n$, and this $\mat{A}$ is the
only matrix that works. \\[0.3em]
\emph{Proof sketch.} Write $\vec{x} = \sum_j x_j \vec{e}_j$. By linearity
$T(\vec{x}) = \sum_j x_j T(\vec{e}_j) = \sum_j x_j (\text{col } j) = \mat{A}\vec{x}$.
Uniqueness: if $\mat{B}$ also represents $T$, then column $j$ of $\mat{B}$ is
$\mat{B}\vec{e}_j = T(\vec{e}_j) = $ column $j$ of $\mat{A}$, so $\mat{B}=\mat{A}$.
\hfill$\square$
\end{formalbox}

That is a lot of payoff for very little work. An object you might think
needs infinitely many instructions -- ``tell me where \emph{every} vector goes''
-- is pinned down by a finite table of numbers: just say where the $n$ basis
vectors go. Linearity fills in all the rest.

\begin{aestheticbox}
A linear map carries infinite information (a rule for the whole continuous plane)
yet is compressed, with zero loss, into a tiny $2\times 2$ box of four numbers.
That compression is the reason linear algebra is computable at all. It is also
the seed of every later idea in this book: change of basis is choosing
\emph{which} vectors to track, eigenvectors are the inputs the machine treats
most simply, and the SVD finds the basis in which the machine looks like pure
stretching. All of it grows from ``store the images of the basis.''
\end{aestheticbox}

\begin{labbox}
Build the matrix from its action. Suppose a linear map $T$ on $\R^2$ satisfies
$T\colvec{1\\0} = \colvec{3\\1}$ and $T\colvec{0\\1} = \colvec{-1\\2}$. \\[0.3em]
(a) Write down $\mat{A}$ by stacking these as columns. \\
(b) Use it to compute $T\colvec{4\\5}$ both by the formula
$4\,T(\vec{e}_1) + 5\,T(\vec{e}_2)$ and by the matrix product $\mat{A}\colvec{4\\5}$,
and check that they agree. \\[0.3em]
\emph{Answer:} $\mat{A} = \begin{bmatrix} 3 & -1 \\ 1 & 2\end{bmatrix}$, and both
routes give $\colvec{7\\14}$.
\end{labbox}

\subsection{Columns and the unit square}

The columns have a direct geometric interpretation. The standard basis vectors $\vec{e}_1,
\vec{e}_2$ are the two sides of the \textbf{unit square} -- the square with
corners $(0,0), (1,0), (1,1), (0,1)$. Their images $\mat{A}\vec{e}_1,
\mat{A}\vec{e}_2$ are the two columns of $\mat{A}$, and they are the two sides of
the \textbf{output parallelogram} the unit square becomes. Drawing the columns as
arrows draws the entire transformation, because the rest of the grid is forced to
follow.

\begin{figure}[H]
\centering
\begin{tikzpicture}[scale=1.0,>=Stealth]
  % input unit square
  \begin{scope}[shift={(0,0)}]
    \draw[->,gray!55] (-0.4,0)--(2.6,0); \draw[->,gray!55] (0,-0.4)--(0,2.6);
    \fill[sourceblue!12] (0,0)--(1,0)--(1,1)--(0,1)--cycle;
    \draw[sourceblue,thick] (0,0)--(1,0)--(1,1)--(0,1)--cycle;
    \draw[->,vecblue,very thick] (0,0)--(1,0) node[below right]{$\vec{e}_1$};
    \draw[->,vecred,very thick] (0,0)--(0,1) node[above left]{$\vec{e}_2$};
    \node at (1.3,-0.9) {input: unit square};
  \end{scope}
  % arrow between
  \draw[->,thick,gray] (3.0,1.0)--(4.2,1.0) node[midway,above]{$\mat{A}$};
  % output parallelogram for A = [[2,1],[0,1]] (a shear+stretch)
  \begin{scope}[shift={(5.0,0)}]
    \draw[->,gray!55] (-0.4,0)--(3.6,0); \draw[->,gray!55] (0,-0.4)--(0,2.6);
    \fill[targetgreen!12] (0,0)--(2,0)--(3,1)--(1,1)--cycle;
    \draw[targetgreen,thick] (0,0)--(2,0)--(3,1)--(1,1)--cycle;
    \draw[->,vecblue,very thick] (0,0)--(2,0) node[below right]{$\mat{A}\vec{e}_1$};
    \draw[->,vecred,very thick] (0,0)--(1,1) node[above left]{$\mat{A}\vec{e}_2$};
    \node at (1.6,-0.9) {output: parallelogram};
  \end{scope}
\end{tikzpicture}
\caption{Reading $\mat{A} = \protect\begin{bmatrix} 2 & 1 \\ 0 & 1\protect\end{bmatrix}$
off the unit square. The columns $\protect\colvec{2\\0}$ and $\protect\colvec{1\\1}$
are the images of $\vec{e}_1$ and $\vec{e}_2$; they form the two sides of the
parallelogram that the square becomes. The whole grid follows.}
\label{ch05:fig-unitsquare}
\end{figure}

\begin{walkbox}
Picture the unit square and its two edge-arrows $\vec{e}_1,\vec{e}_2$. Now imagine
sliding the tip of $\vec{e}_2$ straight up to $(0,3)$ while leaving $\vec{e}_1$
alone. What $2\times 2$ matrix did you just describe, and what happens to the
square's area? Then imagine instead swinging $\vec{e}_2$ down until it points
\emph{along} $\vec{e}_1$. What is the matrix now, and why has the square gone
flat?
\end{walkbox}

\section{The Action on the Grid}
\label{ch05:grid}

Let us make good on the promise that the image of a grid is a grid. Suppose
$\mat{A}\vec{e}_1 = \vec{a}_1$ and $\mat{A}\vec{e}_2 = \vec{a}_2$ are the two
columns. A grid point with integer coordinates $(m,n)$ is the input
$m\vec{e}_1 + n\vec{e}_2$, and by linearity it lands at
\[
  \mat{A}(m\vec{e}_1 + n\vec{e}_2) = m\,\vec{a}_1 + n\,\vec{a}_2.
\]
So the original square lattice $\{m\vec{e}_1 + n\vec{e}_2\}$ becomes the
\emph{sheared} lattice $\{m\vec{a}_1 + n\vec{a}_2\}$ -- still a perfectly regular
grid, just built from the new edge-vectors $\vec{a}_1$ and $\vec{a}_2$ instead of
the old $\vec{e}_1, \vec{e}_2$. Straight lines stay straight (a line is
$\vec{p} + t\vec{d}$, and $\mat{A}(\vec{p}+t\vec{d}) = \mat{A}\vec{p} +
t\,\mat{A}\vec{d}$ is again a line), and parallel lines stay parallel. This is
the visual signature of linearity.

\begin{figure}[H]
\centering
\begin{tikzpicture}[scale=0.62,>=Stealth]
  % left: original grid
  \begin{scope}
    \foreach \i in {-2,...,2}{
      \draw[gray!40] (\i,-2)--(\i,2);
      \draw[gray!40] (-2,\i)--(2,\i);
    }
    \draw[->,vecblue,very thick] (0,0)--(1,0) node[below=1pt,vecblue]{\small $\vec{e}_1$};
    \draw[->,vecred,very thick] (0,0)--(0,1) node[above left=1pt,vecred]{\small $\vec{e}_2$};
    \node at (0,-2.9) {original grid};
  \end{scope}
  % arrow
  \draw[->,thick,gray] (3.1,0)--(4.4,0) node[midway,above]{$\mat{A}$};
  % right: sheared grid for A=[[1,1],[0.5,1]]
  \begin{scope}[shift={(8.4,0)}]
    \foreach \i in {-2,...,2}{
      % lines of constant m: p = i*a1 + n*a2, vary n
      \draw[gray!40] ({\i*1 + (-2)*1},{\i*0.5 + (-2)*1}) -- ({\i*1 + 2*1},{\i*0.5 + 2*1});
      \draw[gray!40] ({(-2)*1 + \i*1},{(-2)*0.5 + \i*1}) -- ({2*1 + \i*1},{2*0.5 + \i*1});
    }
    \draw[->,vecblue,very thick] (0,0)--(1,0.5) node[below right=1pt,vecblue]{\small $\mat{A}\vec{e}_1$};
    \draw[->,vecred,very thick] (0,0)--(1,1) node[above left=1pt,vecred]{\small $\mat{A}\vec{e}_2$};
    \node at (0,-2.9) {image grid};
  \end{scope}
\end{tikzpicture}
\caption{A linear map turns the square grid into a slanted grid. With
$\mat{A}=\protect\begin{bmatrix}1 & 1\\ 0.5 & 1\protect\end{bmatrix}$ the basis
arrows tilt to the new columns $\protect\colvec{1\\0.5}$ and $\protect\colvec{1\\1}$,
and every grid line follows. Lines stay straight and parallels stay parallel --
the grid is sheared, never crumpled.}
\label{ch05:fig-grid}
\end{figure}

\begin{strangbox}
In Strang's course you will spend a lot of time on the equation
$\mat{A}\vec{x} = \vec{b}$: \emph{which} input $\vec{x}$ produces a target output
$\vec{b}$? This chapter is the other half of that conversation. Here we hold
$\mat{A}$ fixed and watch it sweep over \emph{all} inputs at once -- the machine's
view. When Strang draws the ``column picture'' of $\mat{A}\vec{x} = \vec{b}$ as a
combination of columns, he is using exactly our identity $\mat{A}\vec{x} = x_1
\vec{a}_1 + x_2 \vec{a}_2$. Carrying the moving-grid picture in beforehand makes
his row/column pictures feel inevitable instead of arbitrary.
\end{strangbox}

\section{A Tour of $2\times 2$ Machines}
\label{ch05:tour}

We now meet the named transformations of the plane. For each one, the recipe is
the same and it never gets old: \emph{find where $\vec{e}_1$ and $\vec{e}_2$ go,
and stack those as columns.} Keep Figure~\ref{ch05:fig-unitsquare} in mind --
every matrix below is just a different pair of edge-arrows for the unit square.

\subsection{Scalings: stretch and squash}
\label{ch05:scalings}

The simplest machine multiplies each coordinate by a fixed factor. A
\textbf{scaling} with horizontal factor $s$ and vertical factor $t$ sends
$\vec{e}_1 \mapsto \colvec{s\\0}$ and $\vec{e}_2 \mapsto \colvec{0\\t}$, so
\[
  \mat{D} = \begin{bmatrix} s & 0 \\ 0 & t \end{bmatrix},
  \qquad
  \mat{D}\colvec{x\\y} = \colvec{sx\\ty}.
\]
This is a \textbf{diagonal matrix}, $\mat{D} = \diag(s,t)$ -- a matrix whose only
nonzero entries lie along the main diagonal from top-left to bottom-right,
meaning it scales each coordinate axis independently without mixing them. When
$s = t$ it is a uniform zoom (everything grows by the same factor, shapes
preserved). When $s \neq t$ the plane stretches more in one direction than the
other -- a circle becomes an ellipse, as in Figure~\ref{ch05:fig-scaling}. Our
opening example $\diag(2,3)$ was exactly this: double across, triple up.

\begin{recoverybox}
\textbf{Mental check:} What does the matrix $\diag(1,-1)$ do to a vector $\colvec{x\\y}$?
(Pause: it leaves $x$ untouched and flips $y$ to $-y$ -- reflecting every point
across the horizontal axis!)
\end{recoverybox}

\begin{figure}[H]
\centering
\begin{tikzpicture}[scale=1.0,>=Stealth]
  % input unit circle
  \begin{scope}
    \draw[->,gray!55] (-1.4,0)--(1.4,0);
    \draw[->,gray!55] (0,-1.4)--(0,1.5) node[above]{$y$};
    \draw[sourceblue,thick] (0,0) circle (1);
    \draw[->,vecblue,very thick] (0,0)--(1,0) node[below right]{$\vec{e}_1$};
    \draw[->,vecred,very thick] (0,0)--(0,1) node[above left]{$\vec{e}_2$};
    \node at (0,-1.85) {unit circle};
  \end{scope}
  % arrow between
  \draw[->,thick,gray] (1.7,0.3)--(3.6,0.3) node[midway,above=2pt]{$\mat{D}=\diag(2,\tfrac12)$};
  % output ellipse: 2 wide, 1/2 tall
  \begin{scope}[shift={(6.2,0)}]
    \draw[->,gray!55] (-2.3,0)--(2.5,0) node[right]{$x$};
    \draw[->,gray!55] (0,-1.4)--(0,1.5) node[above]{$y$};
    \draw[targetgreen,thick] (0,0) ellipse (2 and 0.5);
    \draw[->,vecblue,very thick] (0,0)--(2,0) node[below right]{$\mat{D}\vec{e}_1$};
    \draw[->,vecred,very thick] (0,0)--(0,0.5) node[above left]{$\mat{D}\vec{e}_2$};
    \node at (0,-1.85) {stretched into an ellipse};
  \end{scope}
\end{tikzpicture}
\caption{A non-uniform scaling $\mat{D}=\diag(2,\tfrac12)$ pulls the unit circle
into an ellipse: $\vec{e}_1$ doubles in length while $\vec{e}_2$ is halved. The
semi-axes of the ellipse are exactly the diagonal factors. The area changes by
the product $2\cdot\tfrac12 = 1$, so this particular stretch keeps area fixed even
though it badly distorts the shape.}
\label{ch05:fig-scaling}
\end{figure}

A negative factor flips that axis: $\diag(-1,1)$ is the reflection across the
$y$-axis from Example~\ref{ch05:three-machines}. A zero factor collapses that axis
entirely: $\diag(1,0)$ flattens the whole plane onto the $x$-axis, which is the
projection $\mat{P}$ from the same example.

\begin{example}[A non-uniform scaling, step by step]
\label{ch05:scaling-ex}
Take $\mat{D} = \diag(\tfrac12, 2)$. The unit square's corner $(1,1)$ goes to
$(\tfrac12, 2)$; the square (area $1$) becomes a rectangle $\tfrac12$ wide and $2$
tall (area $1$ still, since $\tfrac12 \cdot 2 = 1$). The horizontal direction is
squeezed, the vertical stretched, and -- a preview of determinants -- the area is
unchanged because the product of the diagonal factors is $1$.
\end{example}

\subsection{Rotations}
\label{ch05:rotations}

A \textbf{rotation} spins the whole plane rigidly about the origin by an angle
$\theta$ (counterclockwise for $\theta > 0$). Lengths and angles are preserved;
the grid turns like a record on a turntable. To find its matrix, we only need to
know where the two basis arrows swing to.

Spin $\vec{e}_1 = \colvec{1\\0}$, which points east, by $\theta$. Its tip traces
the unit circle, so it lands at $\colvec{\cos\theta\\\sin\theta}$. Spin
$\vec{e}_2 = \colvec{0\\1}$, which points north, by $\theta$; it lands $90^\circ$
ahead of the first arrow, at $\colvec{-\sin\theta\\\cos\theta}$. Stack these as
columns:
\[
  \boxed{\ \mat{R}_\theta
  = \begin{bmatrix} \cos\theta & -\sin\theta \\ \sin\theta & \cos\theta \end{bmatrix}.\ }
\]

\begin{figure}[H]
\centering
\begin{tikzpicture}[scale=2.0,>=Stealth]
  \draw[->,gray!55] (-1.25,0)--(1.35,0) node[right]{$x$};
  \draw[->,gray!55] (0,-0.35)--(0,1.35) node[above]{$y$};
  \draw[gray!35] (0,0) circle (1);
  % e1 and its image
  \draw[->,vecblue,very thick] (0,0)--(1,0) node[below right]{$\vec{e}_1$};
  \draw[->,sourceblue,very thick] (0,0)--({cos(35)},{sin(35)})
       node[above right]{$\mat{R}\vec{e}_1$};
  % e2 and its image
  \draw[->,vecred,very thick] (0,0)--(0,1) node[above right]{$\vec{e}_2$};
  \draw[->,trapred,very thick] (0,0)--({-sin(35)},{cos(35)})
       node[above left]{$\mat{R}\vec{e}_2$};
  % angle arcs showing both basis vectors turn by theta
  \draw[->,targetgreen] (0.45,0) arc (0:35:0.45);
  \node[targetgreen] at (0.62,0.18) {$\theta$};
  \draw[->,targetgreen] (0,0.45) arc (90:125:0.45);
  \node[targetgreen] at (-0.18,0.62) {$\theta$};
\end{tikzpicture}
\caption{Building the rotation matrix from basis images. The east arrow
$\vec{e}_1$ swings to $(\cos\theta,\sin\theta)$ and the north arrow $\vec{e}_2$
swings to $(-\sin\theta,\cos\theta)$. These two images are the columns of
$\mat{R}_\theta$.}
\label{ch05:fig-rotation}
\end{figure}

\begin{example}[Rotation by $90^\circ$]
\label{ch05:rot90}
With $\theta = 90^\circ$ we have $\cos 90^\circ = 0$, $\sin 90^\circ = 1$, so
\[
  \mat{R}_{90^\circ} = \begin{bmatrix} 0 & -1 \\ 1 & 0 \end{bmatrix},
  \qquad
  \mat{R}_{90^\circ}\colvec{a\\b} = \colvec{-b\\a},
\]
matching the rotation we checked in Example~\ref{ch05:check-linear}. The east
arrow $\colvec{1\\0}$ goes to north $\colvec{0\\1}$, and north goes to west
$\colvec{-1\\0}$ -- exactly a quarter turn.
\end{example}

\begin{example}[Rotation by $45^\circ$]
\label{ch05:rot45}
Here $\cos 45^\circ = \sin 45^\circ = \tfrac{1}{\sqrt 2} \approx 0.707$, so
\[
  \mat{R}_{45^\circ} = \frac{1}{\sqrt 2}\begin{bmatrix} 1 & -1 \\ 1 & 1\end{bmatrix}.
\]
Apply it to $\colvec{1\\0}$: out comes $\tfrac{1}{\sqrt2}\colvec{1\\1}$, the
arrow pointing exactly along the diagonal -- a $45^\circ$ tilt, as promised. Note
the columns each have length $1$ and are perpendicular; rotations never change
size or shape, only orientation.
\end{example}

\begin{connectionbox}
The rotation matrices encode the angle-addition formulas of
trigonometry. Rotating by $\alpha$ and then by $\beta$ must be the same as
rotating once by $\alpha + \beta$, i.e. $\mat{R}_\beta \mat{R}_\alpha =
\mat{R}_{\alpha+\beta}$. Multiply the two matrices out and match the top-left
entry:
\[
  \cos(\alpha+\beta) = \cos\beta\cos\alpha - \sin\beta\sin\alpha.
\]
The angle-sum identity you memorized in high school is just the statement that
\emph{composing rotations adds angles}. The same product is how complex
multiplication rotates the plane ($e^{i\alpha} e^{i\beta} = e^{i(\alpha+\beta)}$),
a bridge to the complex-numbers chapter and, later, to the Fourier transform.
\end{connectionbox}

\subsection{Reflections}
\label{ch05:reflections}

A \textbf{reflection} flips the plane across a line through the origin, as if that
line were a mirror. Points on the mirror stay fixed; everything else swaps to the
opposite side. The easiest is reflection across the $x$-axis: it keeps
$\vec{e}_1$ and flips $\vec{e}_2$ to $\colvec{0\\-1}$, giving
\[
  \mat{M}_x = \begin{bmatrix} 1 & 0 \\ 0 & -1 \end{bmatrix}.
\]
More generally, reflection across the line through the origin making angle
$\phi$ with the $x$-axis has matrix
\[
  \mat{M}_\phi = \begin{bmatrix} \cos 2\phi & \sin 2\phi \\ \sin 2\phi & -\cos 2\phi \end{bmatrix}.
\]
(You can derive this by rotating the mirror line down to the $x$-axis, flipping,
and rotating back -- a nice composition exercise once we have the product.) The
$2\phi$ instead of $\phi$ is the famous quirk of reflections: a mirror tilted by
$\phi$ flips directions through twice that angle.

\begin{figure}[H]
\centering
\begin{tikzpicture}[scale=1.5,>=Stealth]
  \draw[->,gray!55] (-1.4,0)--(1.5,0) node[right]{$x$};
  \draw[->,gray!55] (0,-1.4)--(0,1.5) node[above]{$y$};
  % mirror line: x-axis
  \draw[structurepurple,thick,dashed] (-1.3,0)--(1.3,0) node[above left]{mirror};
  % a point and its reflection
  \fill[vecblue] (0.9,0.7) circle (1.3pt) node[above right]{$\vec{x}$};
  \fill[targetgreen] (0.9,-0.7) circle (1.3pt) node[below right]{$\mat{M}_x\vec{x}$};
  \draw[->,vecblue,thick] (0,0)--(0.9,0.7);
  \draw[->,targetgreen,thick] (0,0)--(0.9,-0.7);
  \draw[gray!50,dotted] (0.9,0.7)--(0.9,-0.7);
\end{tikzpicture}
\caption{Reflection across the $x$-axis. The mirror (the $x$-axis) is fixed; the
point $\vec{x}=(0.9,0.7)$ is sent straight down to its mirror image $(0.9,-0.7)$.
The horizontal coordinate is unchanged, the vertical coordinate flips sign.}
\label{ch05:fig-reflection}
\end{figure}

\begin{trapbox}
Rotations and reflections both keep lengths the same, so it is tempting to lump
them together. The difference is \emph{orientation}. A rotation preserves
handedness: a counterclockwise loop stays counterclockwise. A reflection reverses
it: it turns a left hand into a right hand. Algebraically, the determinant of a
rotation is $+1$ and of a reflection is $-1$ -- a sign you will meet again in the
determinants chapter. If you ever confuse the two, ask: does my shape come out
mirror-flipped? If yes, it was a reflection.
\end{trapbox}

\subsection{Shears}
\label{ch05:shears}

A \textbf{shear} slides points parallel to one axis by an amount proportional to
their distance from that axis -- think of pushing the top of a deck of cards
sideways while the bottom stays put. A horizontal shear keeps $\vec{e}_1$ fixed
and tilts $\vec{e}_2$ to $\colvec{k\\1}$:
\[
  \mat{S}_k = \begin{bmatrix} 1 & k \\ 0 & 1 \end{bmatrix},
  \qquad
  \mat{S}_k\colvec{x\\y} = \colvec{x + ky\\y}.
\]
A point's height $y$ is unchanged, but it slides right by $ky$ -- higher points
slide farther. The horizontal axis stays put (when $y = 0$ nothing moves); the rest
of the plane leans like italic type. Notice that a shear does \emph{not} change
area: the parallelogram the unit square becomes has the same base and height as
the square (look back at Figure~\ref{ch05:fig-unitsquare} -- base $1$, height
$1$). This area-preservation is the determinant $\det \mat{S}_k = 1$ again.

\begin{figure}[H]
\centering
\begin{tikzpicture}[scale=1.1,>=Stealth]
  \draw[->,gray!55] (-0.4,0)--(3.4,0) node[right]{$x$};
  \draw[->,gray!55] (0,-0.4)--(0,2.4) node[above]{$y$};
  % original square (faint)
  \draw[sourceblue!50,thick,dashed] (0,0)--(1,0)--(1,1)--(0,1)--cycle;
  % sheared square k=1.2
  \fill[targetgreen!12] (0,0)--(1,0)--(2.2,1)--(1.2,1)--cycle;
  \draw[targetgreen,thick] (0,0)--(1,0)--(2.2,1)--(1.2,1)--cycle;
  \draw[->,vecblue,very thick] (0,0)--(1,0) node[below right]{$\vec{e}_1$ (fixed)};
  \draw[->,vecred,very thick] (0,0)--(1.2,1) node[midway,above left]{$\mat{S}\vec{e}_2$};
  \draw[gray!50,->] (1.1,1.0) -- (2.2,1.0);
  \node[gray, above] at (1.65,1.05) {\small top slides right by $k$};
\end{tikzpicture}
\caption{A horizontal shear $\mat{S}_k$ with $k=1.2$. The base $\vec{e}_1$ is
nailed in place; the top edge slides right in proportion to its height. The dashed
square becomes the green parallelogram -- same base, same height, so the same
area.}
\label{ch05:fig-shear}
\end{figure}

\begin{applicationbox}
These fundamental geometric transformations appear throughout scientific computing.
\emph{Computer graphics and games} build all 2D and 3D motion -- rotating a
character, scaling a sprite, skewing a shadow -- out of exactly these matrices,
applied to the vertices of meshes; a GPU does millions of these matrix--vector
products per frame. \emph{Image processing} uses shears as a practical tool for
image rotation: a rotation can be factored into three shears, and shearing
an image corresponds to shifting rows or columns of pixels in memory.
\emph{Data augmentation} for training neural networks deliberately
applies small random rotations, scalings, reflections, and shears to each training
image so the model learns features that are invariant to these transformations.
\end{applicationbox}

\begin{labbox}
Compute these $2\times 2$ actions by hand and \emph{sketch} the unit square's
image for each. \\[0.3em]
(a) $\mat{R}_{180^\circ} = \begin{bmatrix} -1 & 0 \\ 0 & -1\end{bmatrix}$ on
$\colvec{1\\2}$. \\
(b) The shear $\mat{S}_2 = \begin{bmatrix} 1 & 2 \\ 0 & 1\end{bmatrix}$ on
$\colvec{3\\1}$ and on $\colvec{3\\0}$ -- notice the second is unmoved. Why? \\
(c) The reflection $\mat{M}_x = \begin{bmatrix} 1 & 0 \\ 0 & -1\end{bmatrix}$
applied twice to $\colvec{a\\b}$. What do you get, and what does that say about
$\mat{M}_x\mat{M}_x$? \\[0.3em]
\emph{Answers:} (a) $\colvec{-1\\-2}$. (b) $\colvec{5\\1}$ and $\colvec{3\\0}$;
the second has height $0$, and a shear moves nothing on the axis it fixes. (c)
back to $\colvec{a\\b}$, so reflecting twice is the identity, $\mat{M}_x^2 =
\mat{I}$.
\end{labbox}

\subsection{The identity and the zero map}
\label{ch05:identity-zero}

Two extreme machines bookend the tour. The \textbf{identity}
$\mat{I} = \begin{bmatrix} 1 & 0 \\ 0 & 1\end{bmatrix}$ sends every basis vector
to itself, so $\mat{I}\vec{x} = \vec{x}$: it does nothing, the ``leave it alone''
machine. The \textbf{zero matrix} $\mat{0}$ sends every vector to $\vec{0}$: it
crushes all of space to a single point. Between ``change nothing'' and ``destroy
everything'' live all the interesting transformations.

\section{Composition: Doing One Machine After Another}
\label{ch05:composition}

Machines are made to be chained. Apply $\mat{B}$ to a vector, then apply $\mat{A}$
to the result: ``first $\mat{B}$, then $\mat{A}$.'' The combined effect is itself
a linear map (a composition of linear maps is linear -- chase the definitions),
so by Section~\ref{ch05:columns} it \emph{must} be some matrix. Which matrix?
That matrix is the \textbf{product} $\mat{A}\mat{B}$, defined precisely so that
\[
  \mat{A}(\mat{B}\vec{x}) = (\mat{A}\mat{B})\vec{x} \quad\text{for every } \vec{x}.
\]
This is the real reason matrix multiplication is defined the way it is -- the row-
times-column rule is not arbitrary, it is engineered to make ``compose the
machines'' come out right. We will derive the full rule in the multiplication
chapter; here we just watch it work on the cleanest case and absorb the order.

\begin{example}[Rotate, then rotate: angles add]
\label{ch05:compose-rot}
Compose a $90^\circ$ rotation after another $90^\circ$ rotation. Each is
$\mat{R}_{90^\circ} = \begin{bmatrix} 0 & -1 \\ 1 & 0\end{bmatrix}$, and the
combined machine is found by tracking the basis vectors through both steps. Start
with $\vec{e}_1 = \colvec{1\\0}$:
\[
  \colvec{1\\0}
  \xrightarrow{\ \mat{R}_{90^\circ}\ } \colvec{0\\1}
  \xrightarrow{\ \mat{R}_{90^\circ}\ } \colvec{-1\\0}.
\]
And $\vec{e}_2 = \colvec{0\\1}$:
\[
  \colvec{0\\1}
  \xrightarrow{\ \mat{R}_{90^\circ}\ } \colvec{-1\\0}
  \xrightarrow{\ \mat{R}_{90^\circ}\ } \colvec{0\\-1}.
\]
Stacking the two final images as columns gives the combined matrix
$\begin{bmatrix} -1 & 0 \\ 0 & -1\end{bmatrix} = \mat{R}_{180^\circ}$. Two quarter
turns make a half turn -- exactly what geometry demands, and exactly the
angle-addition we met in the rotation connection box.
\end{example}

\begin{trapbox}
\textbf{Order matters: $\mat{A}\mat{B}$ usually does not equal $\mat{B}\mat{A}$.}
Putting on socks then shoes is not the same as shoes then socks. For a concrete
clash, take a reflection $\mat{M}_x = \begin{bmatrix}1&0\\0&-1\end{bmatrix}$ and a
$90^\circ$ rotation $\mat{R} = \begin{bmatrix}0&-1\\1&0\end{bmatrix}$. Tracking
$\vec{e}_1$: ``rotate then reflect'' sends $\colvec{1\\0}\to\colvec{0\\1}\to
\colvec{0\\-1}$, while ``reflect then rotate'' sends $\colvec{1\\0}\to\colvec{1\\0}
\to\colvec{0\\1}$. Different outputs, so the two orders are different machines.
Also note $\mat{A}\mat{B}$ means ``$\mat{B}$ first'' -- the matrix nearest the
vector acts first. The notation runs right to left, like function composition
$f(g(x))$.
\end{trapbox}

Figure~\ref{ch05:fig-order} draws this clash step by step across two tracks. Tracking where the same arrow $\vec{e}_1$ travels
under both compositions reveals that the two paths end at opposite destinations.

\begin{figure}[H]
\centering
\begin{tikzpicture}[scale=0.85,>=Stealth]
  % Lane 1: M_x * R (Rotate 90deg, then reflect)
  \node[anchor=west,font=\small\bfseries,sourceblue] at (-0.5,1.9) {Order 1: Rotate $90^\circ$, then reflect across $x$-axis ($\mat{M}_x\mat{R}$)};
  
  % Frame 1a: start
  \begin{scope}[shift={(0,0)}]
    \draw[->,gray!50] (-0.8,0)--(1.3,0); \draw[->,gray!50] (0,-0.8)--(0,1.3);
    \draw[->,black,very thick] (0,0)--(1,0) node[below=1pt]{\scriptsize $\vec{e}_1$};
    \node at (0.2,-1.1) {\scriptsize Start: $\colvec{1\\0}$};
  \end{scope}

  \draw[->,very thick,sourceblue] (1.6,0.1) -- (2.8,0.1) node[midway,above]{\scriptsize Rotate $90^\circ$};

  % Frame 1b: rotated
  \begin{scope}[shift={(4.2,0)}]
    \draw[->,gray!50] (-0.8,0)--(1.3,0); \draw[->,gray!50] (0,-0.8)--(0,1.3);
    \draw[->,sourceblue,very thick] (0,0)--(0,1) node[left=1pt]{\scriptsize $\mat{R}\vec{e}_1$};
    \node at (0.2,-1.1) {\scriptsize $\colvec{0\\1}$};
  \end{scope}

  \draw[->,very thick,sourceblue] (5.8,0.1) -- (7.0,0.1) node[midway,above]{\scriptsize Reflect $x$};

  % Frame 1c: reflected
  \begin{scope}[shift={(8.4,0)}]
    \draw[->,gray!50] (-0.8,0)--(1.3,0); \draw[->,gray!50] (0,-1.3)--(0,0.8);
    \draw[->,sourceblue,ultra thick] (0,0)--(0,-1) node[left=1pt]{\scriptsize $\mat{M}_x\mat{R}\vec{e}_1$};
    \node[sourceblue,font=\bfseries] at (0.2,-1.5) {\scriptsize Destination: $\colvec{0\\-1}$};
  \end{scope}

  % Lane 2: R * M_x (Reflect, then rotate 90deg)
  \node[anchor=west,font=\small\bfseries,quotientorange] at (-0.5,-2.5) {Order 2: Reflect across $x$-axis, then rotate $90^\circ$ ($\mat{R}\mat{M}_x$)};

  % Frame 2a: start
  \begin{scope}[shift={(0,-4.4)}]
    \draw[->,gray!50] (-0.8,0)--(1.3,0); \draw[->,gray!50] (0,-0.8)--(0,1.3);
    \draw[->,black,very thick] (0,0)--(1,0) node[below=1pt]{\scriptsize $\vec{e}_1$};
    \node at (0.2,-1.1) {\scriptsize Start: $\colvec{1\\0}$};
  \end{scope}

  \draw[->,very thick,quotientorange] (1.6,-4.3) -- (2.8,-4.3) node[midway,above]{\scriptsize Reflect $x$};

  % Frame 2b: reflected
  \begin{scope}[shift={(4.2,-4.4)}]
    \draw[->,gray!50] (-0.8,0)--(1.3,0); \draw[->,gray!50] (0,-0.8)--(0,1.3);
    \draw[->,quotientorange,very thick] (0,0)--(1,0) node[below=1pt]{\scriptsize $\mat{M}_x\vec{e}_1$};
    \node at (0.2,-1.1) {\scriptsize $\colvec{1\\0}$};
  \end{scope}

  \draw[->,very thick,quotientorange] (5.8,-4.3) -- (7.0,-4.3) node[midway,above]{\scriptsize Rotate $90^\circ$};

  % Frame 2c: rotated
  \begin{scope}[shift={(8.4,-4.4)}]
    \draw[->,gray!50] (-0.8,0)--(1.3,0); \draw[->,gray!50] (0,-0.8)--(0,1.3);
    \draw[->,quotientorange,ultra thick] (0,0)--(0,1) node[left=1pt]{\scriptsize $\mat{R}\mat{M}_x\vec{e}_1$};
    \node[quotientorange,font=\bfseries] at (0.2,-1.1) {\scriptsize Destination: $\colvec{0\\1}$};
  \end{scope}
\end{tikzpicture}
\caption{Matrix multiplication does not commute because the order of physical operations matters.
Top track: rotating $\vec{e}_1$ by $90^\circ$ sends it to $(0,1)$, and then reflecting across the $x$-axis flips it down to $(0,-1)$.
Bottom track: reflecting $\vec{e}_1$ across the $x$-axis leaves it at $(1,0)$, and then rotating by $90^\circ$ swings it up to $(0,1)$.
The two compositions finish at completely opposite destinations: $\mat{M}_x\mat{R} \neq \mat{R}\mat{M}_x$.}
\label{ch05:fig-order}
\end{figure}

\begin{strangbox}
After Strang, you will recognize this section as the doorway to nearly everything
deep. Factoring a matrix is composition read backwards: $\mat{A} = \mat{L}\mat{U}$
says ``$\mat{A}$ is a shear-like machine after another,'' $\mat{A} = \mat{Q}\mat{R}$
splits it into an orthogonal map and a triangular factor, and the singular value
decomposition $\mat{A} = \mat{U}\mathbf{\Sigma}\mat{V}\T$ (the SVD) shows that
\emph{every} matrix is a rotation, then a pure scaling, then another rotation.
Three of the machines from this very chapter -- rotate, stretch, rotate -- compose
to give \emph{any} linear map at all. Hold onto that picture; it is where the book is headed.
\end{strangbox}

\begin{aestheticbox}
Once you read an $m\times n$ matrix as a transformation pinned down by where it sends
the basis vectors, every later topic turns into a geometric question about that motion:
solving $\mat{A}\vec{x}=\vec{b}$ asks which input lands on $\vec{b}$, eigenvectors pick out
the directions the motion leaves unrotated, and the SVD finds perpendicular axes along
which any matrix acts as a pure stretch.
\end{aestheticbox}

\begin{connectionbox}
The area-scaling factor of a linear map is the bridge to multivariable calculus.
A smooth nonlinear map bends the grid, but zoom in close enough to any point and
it looks linear -- its local linear approximation is the \emph{Jacobian matrix} of
partial derivatives. The factor by which that Jacobian scales tiny areas (its
determinant) is precisely the $\det\mat{A}$ you are about to rediscover in
Exercise~\ref{ch05:ex-area}, and it is exactly the $\lvert\partial(x,y)/\partial(u,v)\rvert$
that appears when you change variables in a double integral. ``Matrices as
machines'' is the linearized skeleton of all of differential calculus in several
variables.
\end{connectionbox}

\section{Cost and Intuition Recap}
\label{ch05:recap}

\begin{itemize}
\item \textbf{A matrix is a verb.} An $m\times n$ matrix $\mat{A}$ is the function
  $T(\vec{x})=\mat{A}\vec{x}$ sending $\R^n\to\R^m$. Read the shape first: $n$
  columns in, $m$ rows out (Section~\ref{ch05:matrix-is-function}).
\item \textbf{Linearity is ``the grid stays a grid.''} $T(c\vec{u}+d\vec{v}) =
  cT(\vec{u})+dT(\vec{v})$, and every linear map fixes the origin
  (Section~\ref{ch05:linearity}).
\item \textbf{The columns are the basis images.} Column $j$ is $\mat{A}\vec{e}_j$;
  knowing where the $n$ basis arrows go pins down the entire map
  (Section~\ref{ch05:columns}). Cost to store a linear map on $\R^n\to\R^m$: just
  $mn$ numbers.
\item \textbf{Cost of one action.} Computing $\mat{A}\vec{x}$ for an $m\times n$
  matrix takes $mn$ multiplications and $m(n-1)$ additions -- about $2mn$
  arithmetic operations.
\item \textbf{The named machines.} Scalings $\diag(s,t)$, rotations
  $\mat{R}_\theta$, reflections $\mat{M}_\phi$, shears $\mat{S}_k$ -- each is just
  a choice of two edge-arrows for the unit square (Section~\ref{ch05:tour}).
\item \textbf{Composition is multiplication.} ``First $\mat{B}$, then $\mat{A}$''
  is the matrix $\mat{A}\mat{B}$, and order matters:
  $\mat{A}\mat{B}\neq\mat{B}\mat{A}$ in general (Section~\ref{ch05:composition}).
\end{itemize}

Together, these ideas complete the forward picture: you feed the machine an
input arrow $\vec{x}$, and it carries it to its destination $\mat{A}\vec{x}$.
Yet in real applications, we frequently face the opposite problem: we observe
the destination $\vec{b}$ and must reconstruct the secret input $\vec{x}$ that
arrived there. Running a machine backward -- solving $\mat{A}\vec{x} = \vec{b}$
-- is the central challenge of linear algebra, and the next chapter takes it
head-on.

\begin{historybox}
Before matrices had a name, mathematicians had been computing with determinants
for over a century. A \textbf{determinant} is a single number calculated from a
square grid of numbers that diagnoses whether a set of linear equations has a
unique solution. In 1850, the English mathematician James Joseph Sylvester
(1814--1897) coined the word \textbf{matrix}. Writing in the \emph{Philosophical
Magazine}, Sylvester chose the term from the Latin \emph{matrix}, meaning
``womb'' or ``breeding ground,'' because he viewed the rectangular array of
numbers as a source from which smaller square determinants (minors) could be
extracted. To Sylvester, the determinants were primary and the grid was bookkeeping.

Sylvester's close friend and collaborator Arthur Cayley (1821--1895) treated the
matrix as an algebraic object in its own right. In his 1855 note and his landmark
1858 paper \emph{A Memoir on the Theory of Matrices}, Cayley defined algebraic
operations directly on matrices rather than just on their determinants. He was investigating
\textbf{coordinate substitutions} -- algebraic rules that rename variables to
describe geometric shifts (like replacing $(x,y)$ with $(u,v)$ via
$x = a u + b v$ and $y = c u + d v$). Cayley asked: what happens when you chain
two substitutions together, replacing $(u,v)$ in turn by $(s,t)$ via
$u = \alpha s + \beta t$ and $v = \gamma s + \delta t$? By substituting the
second system into the first and regrouping, Cayley discovered the universal
row-times-column rule for multiplying two matrices.

This new algebra was a significant departure from classical arithmetic. For
centuries, multiplication had been assumed to be \textbf{commutative} ($ab = ba$
-- order never alters the product). But Cayley proved that matrix multiplication
is generally non-commutative ($\mat{A}\mat{B} \neq \mat{B}\mat{A}$), reflecting
the geometric fact that composing spatial transformations in different orders
produces different results. Matrix algebra also admits \textbf{zero divisors}:
two nonzero matrices can multiply to yield the all-zero matrix ($\mat{A}\mat{B}
= \mat{0}$ with $\mat{A}\neq\mat{0}$ and $\mat{B}\neq\mat{0}$). In ordinary arithmetic,
$a \cdot b = 0$ implies $a=0$ or $b=0$; in matrix algebra, two non-trivial transformations
can compose to collapse an entire space to the origin.

Treating transformations as single algebraic objects unlocked massive advances.
In 1925, Werner Heisenberg, Max Born, and Pascual Jordan realized that matrices
were the exact language of quantum mechanics, where physical quantities like
position and momentum fail to commute. Today, matrices form the computational
engine of computer graphics, robotics, and machine learning.

Computational questions about these machines remain open. Cayley's row-times-column
rule multiplies two $n \times n$ matrices using about $n^3$ scalar
multiplications. In 1969, Volker Strassen showed it could be done in $O(n^{2.81})$
operations. Today, the theoretical exponent has been driven down near $2.371$,
but whether two matrices can be multiplied in near-$O(n^2)$ time -- the time
it takes to read the entries -- remains an open problem in computer science.
Beyond multiplication, while running a matrix forward ($\vec{x} \mapsto \mat{A}\vec{x}$)
is effortless arithmetic, running it backward -- finding which input $\vec{x}$
produces a target output $\vec{b}$ -- requires untangling coupled equations.
That inverse challenge takes center stage in the next chapter, which explores
systems of linear equations through the twin perspectives of the row picture
and the column picture.
\end{historybox}

\section{Exercises}
\label{ch05:exercises}

\begin{exercise}[Warmup: act on a vector]
\label{ch05:ex-act}
Let $\mat{A} = \begin{bmatrix} 2 & -1 \\ 1 & 3 \end{bmatrix}$. Compute
$\mat{A}\colvec{1\\0}$, $\mat{A}\colvec{0\\1}$, and $\mat{A}\colvec{1\\1}$. Verify
that the third is the sum of the first two, and explain which property of linear
maps guarantees this.
\end{exercise}

\begin{exercise}[Build the matrix from its action]
\label{ch05:ex-build}
A linear map $T$ on $\R^2$ satisfies $T\colvec{1\\0} = \colvec{2\\-1}$ and
$T\colvec{0\\1} = \colvec{0\\4}$. Write down its matrix, then compute
$T\colvec{3\\2}$ two ways (as a combination of basis images and as a matrix
product) and confirm they agree.
\end{exercise}

\begin{exercise}[Linear or not?]
\label{ch05:ex-linornot}
For each map, decide whether it is linear; if not, give a specific numerical
counterexample to additivity or homogeneity.
(a) $T\colvec{x\\y} = \colvec{x+y\\x-y}$.
(b) $T\colvec{x\\y} = \colvec{x\\y+1}$.
(c) $T\colvec{x\\y} = \colvec{xy\\0}$.
(d) $T\colvec{x\\y} = \colvec{2x\\0}$.
\end{exercise}

\begin{exercise}[Name that machine]
\label{ch05:ex-name}
Identify the geometric effect of each matrix (scaling, rotation, reflection, or
shear), and state where it sends the unit square.
(a) $\begin{bmatrix} 0 & 1 \\ 1 & 0\end{bmatrix}$,
(b) $\begin{bmatrix} 3 & 0 \\ 0 & 3\end{bmatrix}$,
(c) $\begin{bmatrix} 1 & 0 \\ -2 & 1\end{bmatrix}$,
(d) $\begin{bmatrix} 0 & -1 \\ 1 & 0\end{bmatrix}$.
(Hint for (a): where do $\vec{e}_1$ and $\vec{e}_2$ go? What line is the mirror?)
\end{exercise}

\begin{exercise}[Rotation practice]
\label{ch05:ex-rotation}
Write the rotation matrix $\mat{R}_{\theta}$ for $\theta = 30^\circ$ using
$\cos 30^\circ = \tfrac{\sqrt 3}{2}$, $\sin 30^\circ = \tfrac12$. Apply it to
$\colvec{2\\0}$ and verify the output has length $2$ (rotations preserve length).
\end{exercise}

\begin{exercise}[Composition by tracking the basis]
\label{ch05:ex-compose}
Let $\mat{S} = \begin{bmatrix} 1 & 1 \\ 0 & 1\end{bmatrix}$ (shear) and
$\mat{M} = \begin{bmatrix} 1 & 0 \\ 0 & -1\end{bmatrix}$ (reflection). By tracking
$\vec{e}_1$ and $\vec{e}_2$ through both steps, find the matrix for ``first
$\mat{S}$, then $\mat{M}$'' and the matrix for ``first $\mat{M}$, then $\mat{S}$.''
Confirm they are different, illustrating that order matters.
\end{exercise}

\begin{exercise}[Reflection twice]
\label{ch05:ex-involution}
Let $\mat{M}$ be reflection across the $x$-axis. By applying it twice to a general
$\colvec{a\\b}$, show that $\mat{M}^2 = \mat{I}$. A transformation that equals its
own undo button (its own inverse, meaning applying it twice returns every point
to its start: $T^2 = \mat{I}$) is called an \textbf{involution}; name two other
machines from this chapter that are involutions.
\end{exercise}

\begin{exercise}[Collapse and information loss]
\label{ch05:ex-collapse}
Let $\mat{P} = \begin{bmatrix} 1 & 0 \\ 0 & 0\end{bmatrix}$. (a) Where does
$\mat{P}$ send the unit square? (b) Find two \emph{different} input vectors that
$\mat{P}$ sends to the same output, and explain geometrically why $\mat{P}$ cannot
be undone. (This is your first encounter with a non-invertible machine and its
null space.)
\end{exercise}

\begin{exercise}[Proof: linear maps fix the origin]
\label{ch05:ex-origin}
Using only the homogeneity axiom, prove that every linear map satisfies
$T(\vec{0}) = \vec{0}$. Then use additivity to prove $T(-\vec{v}) = -T(\vec{v})$.
Explain in one sentence why this rules out the affine map
$T(\vec{x}) = \vec{x} + \colvec{1\\1}$ being linear.
\end{exercise}

\begin{exercise}[Proof: composition is linear]
\label{ch05:ex-compose-linear}
Let $S, T : \R^2 \to \R^2$ be linear. Prove that their composition $(S \circ
T)(\vec{x}) = S(T(\vec{x}))$ is linear by verifying additivity and homogeneity
directly. (This is why the product of two matrices is again a matrix.)
\end{exercise}

\begin{exercise}[Challenge: derive the general reflection]
\label{ch05:ex-genreflect}
The reflection across the line through the origin at angle $\phi$ can be built as
``rotate the line down to the $x$-axis (by $-\phi$), reflect across the $x$-axis,
rotate back up (by $+\phi$).'' Granting the product rule
$(S\circ T)$'s matrix is the product of the matrices, multiply
$\mat{R}_{\phi}\,\mat{M}_x\,\mat{R}_{-\phi}$ and use double-angle identities to
obtain $\mat{M}_\phi = \begin{bmatrix} \cos 2\phi & \sin 2\phi \\ \sin 2\phi &
-\cos 2\phi\end{bmatrix}$.
\end{exercise}

\begin{exercise}[Challenge: area scaling]
\label{ch05:ex-area}
For each matrix, draw the image of the unit square and compute its area:
$\begin{bmatrix} 2 & 0 \\ 0 & 3\end{bmatrix}$,
$\begin{bmatrix} 1 & 4 \\ 0 & 1\end{bmatrix}$,
$\begin{bmatrix} 2 & 1 \\ 0 & 3\end{bmatrix}$.
Conjecture a formula for the area of the output parallelogram in terms of the
entries $a,b,c,d$ of a general $2\times 2$ matrix $\begin{bmatrix} a & b \\ c & d\end{bmatrix}$
(noting that geometric area is always positive: consider $\abs{ad - bc}$).
You have just rediscovered the determinant!
\end{exercise}
